% !TEX root = ../Thesis.tex
%%
%%  Hochschule für Technik und Wirtschaft Berlin --  Abschlussarbeit
%%
%% Kapitel 4 - Implementierung
%%
%%

\chapter{Implementierung} \label{Implementierung}

\section{Aufbau der Implementierung}

\textbf{Ergänzen}: detaillierte Implementierungsbeschreibung (Hardware-Verdrahtung/Stromversorgung/Schutzgehäuse, Kamera+IR+Polfilter, GPIO/Aktorik-Anbindung). Schaltplan/Blockschaltbild als Abbildung einfügen und im Text referenzieren.

DARSTELLUNG SCHALTUNG

\section{Montage/Vorrichtung}

Im Anwendungsbereich eines Klärwerks ist im Rahmen der kritischen Infrastruktur eine konkrete Planung der benötigten Vorrichtung notwendig, was entsprechend in Absprache mit der Abteilung TS-E/M/Was, deren Spezialisierung in der Montage besteht. Hierfür wurde eine Profilschiene aus Aluminium dahingehend umgebaut, dass eine Fassung für die Reeling des Nachklärbeckens geschaffen wurde, über die sie eingehakt werden kann. Zusätzlich verfügt dieser "Arm" über ein Gelenk, was auch bei eingebauter Vorrichtung die Anpassung des Neigungswinkels ändert und somit indirekt den Kamerawinkel modifiziert. Das ist für Testzwecke notwendig um den Winkel mit der geringsten Spiegelung für das Kameramodul zu ermitteln. \\ \\

\textbf{Technische Zeichnung}/Foto der Vorrichtung ergänzen (Maße, Material, Befestigungspunkte, Neigungswinkelbereich). Zusätzlich: Witterungsschutz, Kabelmanagement, Zugriff für Wartung/Reinigung beschreiben.
TECHNISCHE ZEICHNUNG \\ \\

PROTOKOLL \\ \\